\documentclass[10pt, letterpaper]{article}

% Packages:
\usepackage[
    ignoreheadfoot, % set margins without considering header and footer
    top=2 cm, % seperation between body and page edge from the top
    bottom=2 cm, % seperation between body and page edge from the bottom
    left=2 cm, % seperation between body and page edge from the left
    right=2 cm, % seperation between body and page edge from the right
    footskip=1.0 cm, % seperation between body and footer
    % showframe % for debugging 
]{geometry} % for adjusting page geometry
\usepackage{titlesec} % for customizing section titles
\usepackage{tabularx} % for making tables with fixed width columns
\usepackage{array} % tabularx requires this
\usepackage[dvipsnames]{xcolor} % for coloring text
\definecolor{primaryColor}{RGB}{0, 0, 0} % define primary color
\usepackage{enumitem} % for customizing lists
\usepackage{fontawesome5} % for using icons
\usepackage{amsmath} % for math
\usepackage[
    pdftitle={DavidChan-Resume},
    pdfauthor={David Chan},
    pdfcreator={LaTeX with RenderCV},
    colorlinks=true,
    urlcolor=primaryColor
]{hyperref} % for links, metadata and bookmarks
\usepackage[pscoord]{eso-pic} % for floating text on the page
\usepackage{calc} % for calculating lengths
\usepackage{bookmark} % for bookmarks
\usepackage{lastpage} % for getting the total number of pages
\usepackage{changepage} % for one column entries (adjustwidth environment)
\usepackage{paracol} % for two and three column entries
\usepackage{ifthen} % for conditional statements
\usepackage{needspace} % for avoiding page brake right after the section title
\usepackage{iftex} % check if engine is pdflatex, xetex or luatex

% Ensure that generate pdf is machine readable/ATS parsable:
\ifPDFTeX
    \input{glyphtounicode}
    \pdfgentounicode=1
    \usepackage[T1]{fontenc}
    \usepackage[utf8]{inputenc}
    \usepackage{lmodern}
\fi

\usepackage{charter}

% Some settings:
\raggedright
\AtBeginEnvironment{adjustwidth}{\partopsep0pt} % remove space before adjustwidth environment
\pagestyle{empty} % no header or footer
\setcounter{secnumdepth}{0} % no section numbering
\setlength{\parindent}{0pt} % no indentation
\setlength{\topskip}{0pt} % no top skip
\setlength{\columnsep}{0.15cm} % set column seperation
\pagenumbering{gobble} % no page numbering

\titleformat{\section}{\needspace{4\baselineskip}\bfseries\large}{}{0pt}{}[\vspace{1pt}\titlerule]

\titlespacing{\section}{
    % left space:
    -1pt
}{
    % top space:
    0.3 cm
}{
    % bottom space:
    0.2 cm
} % section title spacing

\renewcommand\labelitemi{$\vcenter{\hbox{\small$\bullet$}}$} % custom bullet points
\newenvironment{highlights}{
    \begin{itemize}[
        label=-,
        topsep=0.10 cm,
        parsep=0.10 cm,
        partopsep=0pt,
        itemsep=0pt,
        leftmargin=0 cm + 10pt
    ]
}{
    \end{itemize}
} % new environment for highlights


\newenvironment{highlightsforbulletentries}{
    \begin{itemize}[
        topsep=0.10 cm,
        parsep=0.10 cm,
        partopsep=0pt,
        itemsep=0pt,
        leftmargin=10pt
    ]
}{
    \end{itemize}
} % new environment for highlights for bullet entries

\newenvironment{onecolentry}{
    \begin{adjustwidth}{
        0 cm + 0.00001 cm
    }{
        0 cm + 0.00001 cm
    }
}{
    \end{adjustwidth}
} % new environment for one column entries

\newenvironment{twocolentry}[2][]{
    \onecolentry
    \def\secondColumn{#2}
    \setcolumnwidth{\fill, 4.5 cm}
    \begin{paracol}{2}
}{
    \switchcolumn \raggedleft \secondColumn
    \end{paracol}
    \endonecolentry
} % new environment for two column entries

\newenvironment{threecolentry}[3][]{
    \onecolentry
    \def\thirdColumn{#3}
    \setcolumnwidth{, \fill, 4.5 cm}
    \begin{paracol}{3}
    {\raggedright #2} \switchcolumn
}{
    \switchcolumn \raggedleft \thirdColumn
    \end{paracol}
    \endonecolentry
} % new environment for three column entries

\newenvironment{header}{
    \setlength{\topsep}{0pt}\par\kern\topsep\centering\linespread{1.5}
}{
    \par\kern\topsep
} % new environment for the header

\newcommand{\placelastupdatedtext}{% \placetextbox{<horizontal pos>}{<vertical pos>}{<stuff>}
  \AddToShipoutPictureFG*{% Add <stuff> to current page foreground
    \put(
        \LenToUnit{\paperwidth-2 cm-0 cm+0.05cm},
        \LenToUnit{\paperheight-1.0 cm}
    ){\vtop{{\null}\makebox[0pt][c]{
        \small\color{gray}\textit{Last updated in September 2024}\hspace{\widthof{Last updated in September 2024}}
    }}}%
  }%
}%

% save the original href command in a new command:
\let\hrefWithoutArrow\href

% new command for external links:

%==================================================

\begin{document}
    \newcommand{\AND}{\unskip
        \cleaders\copy\ANDbox\hskip\wd\ANDbox
        \ignorespaces
    }
    \newsavebox\ANDbox
    \sbox\ANDbox{$|$}

    \begin{header}

        \fontsize{25 pt}{25 pt}\selectfont \textbf{David Le Chan}
        \vspace{1 pt}

        \normalsize
        \mbox{Los Altos, CA}%
        \kern 3.0 pt%
        \AND%
        \kern 3.0 pt%
        \mbox{\hrefWithoutArrow{mailto:davidlechan@gmail.com}{davidlechan@gmail.com}}%
        \kern 3.0 pt%
        \AND%
        \kern 3.0 pt%
        \mbox{\hrefWithoutArrow{tel:+1-650-383-8954}{(650) 383-8954}} 
        
        \mbox{\hrefWithoutArrow{https://davidlechan.dev/}{\underline{davidlechan.dev}}}%
        \kern 3.0 pt%
        \AND%/
        \kern 3.0 pt%
        \mbox{\href{https://linkedin.com/in/david-le-chan}{\underline{linkedin.com/in/david-le-chan}}}%
        \kern 3.0 pt%
        \AND%
        \kern 3.0 pt%
        \mbox{\hrefWithoutArrow{https://github.com/code49}{\underline{github.com/code49}}}%
    
    \end{header}

    % \vspace{5 pt - 0.3 cm} %==================================================

    \section{Education}

        \begin{twocolentry}
            {Expected May 2027}
            {\textbf{M.S. in Electrical and Computer Engineering (ECE)} \\ 
            \href{https://www.ece.cmu.edu/academics/ms-ece/index.html}{\underline{Carnegie Mellon University (CMU)}} | Pittsburgh, PA}
        \end{twocolentry}
        \vspace{0.10 cm}
        \begin{onecolentry}
            \begin{highlights}
                % \item GPA: /4.00; \textit{x College of Engineering Dean's List Recipient}
                \item Reconfigurable Logic, Hardware Verification, Digital Integrated Circuit Design (SRAMs), %
                Post-Silicon Hardware Testing, Memory Devices, Innovation Strategy and Management
            \end{highlights}
        \end{onecolentry}

        \vspace{0.2 cm} %==================================================
    
        \begin{twocolentry}
            {Expected May 2026}
            {\textbf{B.S. in Electrical and Computer Engineering (ECE)} \\ 
            \href{https://www.ece.cmu.edu/academics/bs-in-ece/index.html}{\underline{Carnegie Mellon University (CMU)}} | Pittsburgh, PA}
        \end{twocolentry}
        \vspace{0.10 cm}
        \begin{onecolentry}
            \begin{highlights}
                \item GPA: 3.8/4.0; \textit{6x College of Engineering Dean's List Recipient}
                \item Digital IC Tapeout, Computer Architecture, Digital Design Verification, %
                Microelectronic Circuits, Machine Learning, Computer Systems, Signal Analysis, %
                Linear Algebra, Multivariable Calculus, Statistics, Leadership
            \end{highlights}
        \end{onecolentry}

    \section{Work Experience}

        \begin{twocolentry} 
            {May 2025 - August 2025}
            {\textbf{FPGA \& Electrical Engineering Intern} \\
            \href{https://www.kla.com/}{\underline{KLA Corporation}, BBP Division} | Milpitas, CA}
        \end{twocolentry}
        \vspace{0.10 cm}
        \begin{onecolentry}
            \begin{highlights}
                \item Upgraded FPGA firmware for high-speed lossless image compression on KLA's flagship wafer inspection tool line
                \item Implemented subsystems via Verilog and Vivado IP Integrator, verifying functionality through Questa simulation
                \item Deployed designs on PCIe-based Alveo accelerator cards, performing place-and-route optimization and hardware-level validation to quantify performance and ensure reliability
            \end{highlights}
        \end{onecolentry}

        \vspace{0.2 cm} %==================================================

        \begin{twocolentry}
            {May 2024 - Present}
            {\textbf{Undergraduate Research Assistant} \\
            IO Harness Project, CMU ECE | Pittsburgh, PA}
        \end{twocolentry}
        \vspace{0.10 cm}
        \begin{onecolentry}
            \begin{highlights}

                % \item Designing a standardized chip harness on TSMC’s 180nm node with Professors Ken Mai and Jim Bain to reduce infrastructure redevelopment work in CMU’s digital IC tapeout (18-725) class
                \item Designing a standardized harness to reduce infrastructure redevelopment work for IC tapeouts; first prototype fabricated on TSMC 180\,nm under active testing
                \item Architecting system features and writing SystemVerilog RTL for I2C, UART, and SPI communication blocks

            \end{highlights}
        \end{onecolentry}

        \vspace{0.2 cm} %==================================================

        \begin{twocolentry}
            {June 2022 - January 2024}
            {\textbf{Power Electronics \& Programming Intern} \\
            \href{https://www.taumotors.com/}{\underline{Tau Motors}} | Redwood City, CA}
        \end{twocolentry}
        \vspace{0.10 cm}
        \begin{onecolentry}
            \begin{highlights}

                \item Prototyped power circuits for wound-field electric motors, including PCB layout, assembly, and bench testing
                \item Developed custom Python-based inventory management software and systems to accelerate hardware iterations

            \end{highlights}
        \end{onecolentry}

    % \section{Publications}
        
    %     \begin{samepage}
    %         \begin{twocolentry}{
    %             Jan 2004
    %         }
    %             \textbf{3D Finite Element Analysis of No-Insulation Coils}
    %         \end{twocolentry}

    %         \vspace{0.10 cm}
            
    %         \begin{onecolentry}
    %             \mbox{Frodo Baggins}, \mbox{\textbf{\textit{John Doe}}}, \mbox{Samwise Gamgee}

    %             \vspace{0.10 cm}
                
    %     \href{https://doi.org/10.1109/TASC.2023.3340648}{10.1109/TASC.2023.3340648}
    %     \end{onecolentry}
    %     \end{samepage}
    
    \section{Leadership and Projects}

        \begin{twocolentry}
            {January 2025 - January 2026}
            {\textbf{Head Teaching Assistant (TA)}, Introduction to ECE (18-100) \\}
        \end{twocolentry}
        \vspace{0.10 cm}
        \begin{onecolentry}
            \begin{highlights}

                \item Leading a 40 TA team to foster an emotionally-safe environment where 180 first-year students can explore ECE, form lasting friendships, and develop strong engineering habits
                \item Continuously redesigning lab curriculum, such as the AM Radio lab with PCB and soldering components, to give students hands-on experience with real hardware assembly and industry design workflows
                \item Establishing automated Python scripts and feedback systems to streamline course operations, reducing administrative overhead and empowering TAs to focus on mentorship and teaching quality

            \end{highlights}
        \end{onecolentry}

        \vspace{0.2 cm} %==================================================
        
        \begin{twocolentry}
            {January 2025}
            \textbf{One-Instruction Flappy Bird }\href{https://github.com/jobitaki/JustOneFlappyBird}{\underline{(github.com/jobitaki/JustOneFlappyBird)}}
        \end{twocolentry}
        \vspace{0.10 cm}
        \begin{onecolentry}
            \begin{highlights}
                \item Collaborated with a team to design a one-instruction (SUBLEQ) CPU to play the video game Flappy Bird
                % \item Collaborated with a team to design a one-instruction (SUBLEQ) CPU to play the video game Flappy Bird, contributing to a successful project at CMU ECE’s Build18 hackathon
                \item Implemented memory-mapped IO and VGA graphics features in SystemVerilog to target a Spartan7 FPGA
            \end{highlights}
        \end{onecolentry}


        \vspace{0.2 cm} %==================================================

        % \begin{twocolentry}
        %     {August 2023 - Present}
        %     \href{https://tartanconnect.cmu.edu/cia/top/}{\textbf{Carnegie Involvement Association (CIA) Buggy}}
        % \end{twocolentry}
        % \vspace{0.10 cm}
        % \begin{onecolentry}
        %     \begin{highlights}
        %         \item Managed over \$50k in club funds and coordinated with CMU alumni and financial officers
        %         % \item Pusher and mechanic for buggy, CMU's annual unpowered vehicle race, as part of team CIA
        %     \end{highlights}
        % \end{onecolentry}

        % \vspace{0.2 cm} %==================================================
    
    % \section{Skills}

    %     \begin{onecolentry}
    %         \begin{highlights}
    %             \item SystemVerilog, VCS/Questa, Quartus/Vivado, TCL, Cadence (Virtuoso/Genus/Innovus), KiCAD/Fusion360
    %             \item Python (NumPy, Pandas, Scikit), Git, Bash, C/C++, MATLAB
    %             % \item English (Native), Mandarin Chinese (Proficient), Git, LaTeX, Linux, MS Office
    %         \end{highlights}
    %     \end{onecolentry}

\end{document}

%==================================================
 
% TEMPLATE ITEM

% \begin{twocolentry} 
%     {begin - end}
%     {\textbf{title} \\
%     \href{site_url}{company} | location}
% \end{twocolentry}
% \vspace{0.10 cm}
% \begin{onecolentry}
%     \begin{highlights}
%
%         \item bullet1
%         \item bullet2
%
%     \end{highlights}
% \end{onecolentry}

%==================================================
